
\begin{sublemme}
\label{XIII.6.1.2}
Sean
\marginpar{438}
$S$ un esquema localmente noetheriano y $F\to G$ un morfismo
sobreyectivo de haces
\ifthenelse{\boolean{orig}}
{de grupos}
{de conjuntos}
sobre~$S$. Sean los $q_i$ una familia finita de secciones de~$G$
\ifthenelse{\boolean{orig}}
{sobre~$X$}
{sobre~$S$}
que engendran~$G$, y, para cada~$i$, sea $F_i$ el subhaz de~$F$ que es
la imagen inversa de~$q_i$. Entonces, si $G$ y los $F_i$ son
constructibles, también $F$ es constructible.
\end{sublemme}

Para demostrar que $F$ es constructible, basta probar que, para todo
punto~$s$ de~$S$, existe un entorno abierto~$U$ de~$s$ tal que $F|U$
sea localmente constante y constructible. Sea, pues, $s$ un punto de~$S$.
Como los haces $F_i$ y $G$ son constructibles, puede encontrarse un
entorno abierto~$U$ de~$s$ tal que $F_i|U$ y $G|U$ sean localmente
constantes. Demostremos entonces que $F|U$ es localmente constante.
Según (SGA~4 IX~2.13~(i)), basta ver que, si $\overline{s}$ es un punto
geométrico situado sobre~$s$, $\tilde{s}$ un punto geométrico de~$U$ y
$\overline{s}\to\tilde{s}$ un morfismo de especialización, el morfismo
canónico
$$
F_{\tilde{s}}\to F_{\overline{s}}
$$
es biyectivo.

Consideremos los diagramas conmutativos
$$
\xymatrix{ (F_i)_{\tilde{s}} \ar[r]^-{\tilde{a}_i} \ar[d]^{\wr} &
F_{\tilde{s}} \ar[r]^-{\tilde{a}} \ar[d]_b & G_{\tilde{s}}
\ar[d]^{\wr}\\ (F_i)_{\overline{s}} \ar[r]^-{\overline{a}_i} &
F_{\overline{s}} \ar[r]^-{\overline{a}} & G_{\overline{s}} }
$$
Sea $\overline{q}_i$ (\resp $\tilde{q}_i$) la imagen inversa de~$q_i$
en~$G_{\overline{s}}$ (\resp en~$G_{\tilde{s}}$); los morfismos
$\overline{a}$ y $\tilde{a}$ son sobreyectivos, y el morfismo
$\overline{a}_i$ (\resp $\tilde{a}_i$) induce una biyección de
$(F_i)_{\overline{s}}$ (\resp $(F_i)_{\tilde{s}}$) sobre
\ifthenelse{\boolean{orig}}
{$\overline{a}^{-1}(q_i)$ (\resp sobre~$\tilde{a}^{-1}(q_i)$)}
{$\overline{a}^{-1}(\overline{q}_i)$
(\resp sobre~$\tilde{a}^{-1}(\tilde{q}_i)$)}. Se deduce,
pues, del diagrama anterior que $b$ es un isomorfismo.
